\chapter{Referencial Teórico}
O objetivo desta seção é fundamentar teoricamente a pesquisa, delimitando os conceitos, princípios e 
parâmetros institucionais que norteiam o estudo. Este capítulo visa conferir rigor metodológico e 
consistência analítica à avaliação da gestão pública municipal e às suas ferramentas de governança.

\section{Índice de Efetividade da Gestão Municipal}
Concebido originalmente em 2014 pelo Tribunal de Contas do Estado de São Paulo (TCE-SP), o Índice de 
Efetividade da Gestão Municipal (IEGM) foi posteriormente incorporado e aperfeiçoado pelo Instituto 
Rui Barbosa (IRB) - entidade voltada ao desenvolvimento de estudos e pesquisas para o sistema de 
controle externo \cite{iegm_ceara_2022}. Com o intuito de uniformizar a avaliação das administrações 
públicas locais, o índice foi selecionado em 2016 pela Rede Nacional de Indicadores Públicos 
(Indicon) para aplicação em âmbito nacional \cite{iegm_ceara_2022}.

O indicador foi projetado para mensurar o nível de eficácia e desenvolvimento das 
administrações municipais brasileiras. O IEGM engloba sete áreas fundamentais: educação, saúde, 
gestão fiscal, planejamento, meio ambiente, defesa civil e governança de TI \cite{lima_influencia_2023}.

O resultado geral do IEGM varia de 0 a 100, no qual a pontuação máxima representa a aderência 
plena do município às normas e aos controles estabelecidos. Para o cálculo do índice consolidado, 
aplica-se uma métrica de ponderação que reflete a relevância de cada dimensão específica. Assim, a 
composição do IEGM é estruturada pelas seguintes proporções:
\cite{lima_influencia_2023}.

\begin{itemize}
  \item Educação (i-Educ, 20\%)
  \item Saúde (i-Saúde, 20\%)
  \item Gestão Fiscal (i-Fiscal, 20\%)
  \item Meio Ambiente (i-Amb, 10\%)
  \item Proteção dos Cidadãos/Defesa Civil (i-Cidade, 5\%)
  \item Governança de TI (i-Gov TI, 5\%)
\end{itemize}

O Índice conta com a adesão de todos os 27 Tribunais de Contas dos estados e do Distrito Federal, 
consolidando-se como um instrumento nacional de avaliação das administrações públicas locais \cite{araujo_gestoes_municipais_2021}. 
A operacionalização do índice ocorre por meio do Termo de Adesão à Rede de Indicadores Nacionais, 
firmado entre o Instituto Rui Barbosa (IRB) e os órgãos de contas estaduais e distrital. A partir 
desse acordo, as informações prestadas pelos municípios são coletadas e posteriormente validadas 
in loco ou auditadas pelos respectivos tribunais, garantindo a consistência dos dados do índice \cite{araujo_gestoes_municipais_2021}.

\section{Governança de TI na Administração Pública}

\section{Governança Frugal de TI em Municípios de Baixa Maturidade}

\section{Design Science Research(DSR) como Abordagem Metodológica}

\section{Trabalhos Relacionados}


